% META: {"id":"TNSI-LANGAGES-ET-PROGRAMMATION-RE-C05","chapitre":"TNSI-LANGAGES-ET-PROGRAMMATION","type_objet":"remediation","capacites_codes":["C11"],"status":"approved","origin":"generated","status_history":["generated","audited","accepted","approved"],"release_acceptance":"RELEASE_OWNER_BATCH_ACCEPTANCE","acceptance_closure_digest":"sha256:9b3ccf9a81c5520fbb7e03b7b2d3e7bf2057a02834b6d908bf3be5f49f3b553f","acceptance_receipt_digest":"sha256:4fad86f0282e0fa4e0419cca1ad6379ae7af5a280c04a4a90880f90a1c044242"}

%---------------------------------------------------------------
% FICHE DE REMEDIATION — C05 : inégalité stricte/large (off-by-one)
%---------------------------------------------------------------

\begin{exercice}{TNSI-LANG-RE-C05-EX1}{1}{10}
Une fonction censée vérifier qu'un indice est valide pour une liste de longueur \lstinline{n} est écrite ainsi :
\begin{python}
def indice_valide_bug(indice, n):
    return 0 <= indice <= n
\end{python}
Pour une liste de longueur $n = 5$ (indices valides : $0$ à $4$), que renvoie \lstinline{indice_valide_bug(5, 5)} ? Est-ce correct ? Corriger la fonction.
\end{exercice}

% BEGIN-VERIFY
% def indice_valide_bug(indice, n):
%     return 0 <= indice <= n
% assert indice_valide_bug(5, 5) is True
% def indice_valide_corrige(indice, n):
%     return 0 <= indice < n
% assert indice_valide_corrige(5, 5) is False
% assert indice_valide_corrige(4, 5) is True
% END-VERIFY

\begin{corrige}{TNSI-LANG-RE-C05-EX1}
\lstinline{indice_valide_bug(5, 5)} renvoie \lstinline{True}, ce qui est \textbf{incorrect} : pour une liste de longueur $5$, les indices valides sont $0, 1, 2, 3, 4$ (indice $5$ est hors limites, il provoquerait un \lstinline{IndexError}). L'erreur vient de l'inégalité \textbf{large} \lstinline{<= n} au lieu de \textbf{stricte} \lstinline{< n} :
\begin{python}
def indice_valide_corrige(indice, n):
    return 0 <= indice < n
\end{python}
\end{corrige}
